\documentclass[11pt,a4paper]{article}
\usepackage{fontspec,amsmath,amssymb,amsthm,microtype}
\setmainfont{lmroman10-regular.otf}[BoldFont=lmroman10-bold.otf,ItalicFont=lmroman10-italic.otf,BoldItalicFont=lmroman10-bolditalic.otf]
\usepackage{xcolor}
\usepackage[most]{tcolorbox}
\usepackage{mathrsfs}
\providecommand{\lt}{<}
\providecommand{\gt}{>}
\setmonofont{DejaVuSansMono.ttf}[BoldFont=DejaVuSansMono-Bold.ttf,ItalicFont=DejaVuSansMono-Oblique.ttf,BoldItalicFont=DejaVuSansMono-BoldOblique.ttf,Scale=0.86]
\newfontfamily\leanfallback{STIXTwoMath-Regular.otf}[Scale=0.86]
\definecolor{leanrule}{gray}{0.62}
\definecolor{leanback}{gray}{0.962}
\definecolor{leantitle}{gray}{0.22}
\definecolor{leanmuted}{gray}{0.38}
\newtcolorbox{leanbox}[1]{enhanced,breakable,colback=leanback,colframe=leanrule,
  boxrule=0.4pt,arc=0.6mm,left=2.4mm,right=2.4mm,top=1.8mm,bottom=1.8mm,
  before skip=0.8em,after skip=1em,fonttitle=\small\bfseries,coltitle=leantitle,
  colbacktitle=leanback,titlerule=0.3pt,title={#1}}
\newlength{\leancharwidth}
\newenvironment{leancode}{\par\smallskip\begingroup\ttfamily\footnotesize\raggedright
  \settowidth{\leancharwidth}{M}\setlength{\parskip}{0pt}\setlength{\parindent}{0pt}
  \frenchspacing}{\par\endgroup\smallskip}
\newcommand{\leanline}[2]{\hangindent=\dimexpr#1\leancharwidth+4\leancharwidth\relax
  \hangafter=1\noindent\hspace*{#1\leancharwidth}#2\par}
\newcommand{\leansource}[1]{\par{\footnotesize\color{leanmuted}#1}\par}
\newcommand{\leanlabel}[1]{\par\smallskip\noindent{\small\bfseries #1}\par}
\usepackage[margin=25mm]{geometry}
\usepackage[unicode,hidelinks]{hyperref}
\hypersetup{pdfauthor={LangevinLean Project}}
\setlength{\emergencystretch}{3em}
\setlength{\parskip}{0.4em}
\setlength{\parindent}{0pt}
\allowdisplaybreaks[2]
\raggedbottom
\title{\texorpdfstring{Proximal points and the subgradient optimality condition}{Proximal points and the subgradient optimality condition}}
\author{LangevinLean Project}
\date{}
\begin{document}\maketitle
Draft mathematical exposition; these notes may contain mistakes.


Throughout this section \(E\) is a real inner-product space with inner product \(\langle\cdot,\cdot\rangle\) and norm \(\|\cdot\|\), \(f:E\to\mathbb R\) is a real-valued function and \(\lambda>0\). The \emph{proximal objective} at an input \(x\in E\) is \(Q_{\lambda,x}(p)=f(p)+\|p-x\|^2/(2\lambda)\), and a \emph{proximal point} of \(f\) at \(x\) is a global minimizer of \(Q_{\lambda,x}\). A \emph{subgradient} of \(f\) at \(p\) is a vector \(s\) such that the affine function \(y\mapsto f(p)+\langle s,y-p\rangle\) lies below \(f\) on all of \(E\).




We prove: (i) subgradients are monotone; (ii) \(p\) is a proximal point at \(x\) whenever the scaled residual \((x-p)/\lambda\) is a subgradient of \(f\) at \(p\), and, for convex \(f\), conversely; (iii) for convex \(f\), proximal points satisfy the firm nonexpansiveness inequality \(\|p-q\|^2\le\langle x-y,p-q\rangle\), so they are unique and both \(x\mapsto p\) and \(x\mapsto x-p\) are \(1\)-Lipschitz; (iv) if \(E\) is finite-dimensional and \(f\) is convex and bounded below, a proximal point exists at every input. The proximal map \(P_\lambda\) is then the map sending \(x\) to its unique proximal point.




These facts are used for the Moreau envelope in the next section, whose gradient is \((x-P_\lambda(x))/\lambda\), for the energy inequalities of the pairwise proximal particle updates, and for the contraction and parameter-sensitivity estimates of the splitting scheme.


\section{\texorpdfstring{Proximal objective, proximal points and monotone subgradients}{Proximal objective, proximal points and monotone subgradients}}


\textbf{Definitions.} Let \(f:E\to\mathbb R\), \(\lambda>0\) and \(x\in E\). The proximal objective is


\[Q_{\lambda,x}(p)=f(p)+\frac{\|p-x\|^2}{2\lambda}\qquad(p\in E).\]


A point \(p\in E\) is a \emph{proximal point of \(f\) at \(x\) with parameter \(\lambda\)} if \(Q_{\lambda,x}(p)\le Q_{\lambda,x}(y)\) for every \(y\in E\), that is, if \(p\) is a global minimizer of \(Q_{\lambda,x}\). A vector \(s\in E\) is a \emph{subgradient of \(f\) at \(p\)}, written \(s\in\partial f(p)\), if


\[f(p)+\langle s,y-p\rangle\le f(y)\qquad\text{for every }y\in E.\qquad(1)\]


Inequality (1) says that the affine function \(y\mapsto f(p)+\langle s,y-p\rangle\) lies below \(f\) on all of \(E\) and agrees with \(f\) at \(p\). Neither definition requires convexity or differentiability of \(f\).




\textbf{Claim (monotonicity of subgradients).} Let \(f:E\to\mathbb R\) be any function, let \(p,q\in E\), \(s\in\partial f(p)\) and \(t\in\partial f(q)\). Then


\[\langle s-t,p-q\rangle\ge0.\]


\textbf{Proof.} Apply (1) for \(s\) at \(p\) with \(y=q\), and (1) for \(t\) at \(q\) with \(y=p\):


\[f(p)+\langle s,q-p\rangle\le f(q),\qquad f(q)+\langle t,p-q\rangle\le f(p).\]


Since \(q-p=-(p-q)\), linearity of the inner product in its second argument gives \(\langle s,q-p\rangle=-\langle s,p-q\rangle\). Adding the two inequalities and cancelling the real number \(f(p)+f(q)\) from both sides gives


\[-\langle s,p-q\rangle+\langle t,p-q\rangle\le0.\]


By linearity in the first argument the left-hand side equals \(-\langle s-t,p-q\rangle\). Hence \(\langle s-t,p-q\rangle\ge0\). \(\square\)

\begin{leanbox}{Lean: \texttt{proxObjective}, \texttt{IsProxPoint}, \texttt{IsSubgradient} and 1 more}
\leansource{\texttt{proxObjective} \quad ProximalLangevin/\allowbreak{}Proximal.lean, lines 30–31}
\begin{leancode}
\leanline{0}{def\ proxObjective\ (f\ :\ E\ →\ ℝ)\ (lam\ :\ ℝ)\ (x\ p\ :\ E)\ :\ ℝ\ :=}
\leanline{2}{f\ p\ +\ ‖p\ -\ x‖\ \textasciicircum{}\ 2\ /\ (2\ *\ lam)}
\end{leancode}
\leansource{\texttt{IsProxPoint} \quad ProximalLangevin/\allowbreak{}Proximal.lean, lines 34–35}
\begin{leancode}
\leanline{0}{def\ IsProxPoint\ (f\ :\ E\ →\ ℝ)\ (lam\ :\ ℝ)\ (x\ p\ :\ E)\ :\ Prop\ :=}
\leanline{2}{∀\ y,\ proxObjective\ f\ lam\ x\ p\ ≤\ proxObjective\ f\ lam\ x\ y}
\end{leancode}
\leansource{\texttt{IsSubgradient} \quad ProximalLangevin/\allowbreak{}Proximal.lean, lines 38–39}
\begin{leancode}
\leanline{0}{def\ IsSubgradient\ (f\ :\ E\ →\ ℝ)\ (p\ g\ :\ E)\ :\ Prop\ :=}
\leanline{2}{∀\ y,\ f\ p\ +\ {\leanfallback ⟪}g,\ y\ -\ p{\leanfallback ⟫}\ ≤\ f\ y}
\end{leancode}
\leansource{\texttt{IsSubgradient.\allowbreak{}monotone} \quad ProximalLangevin/\allowbreak{}Proximal.lean, lines 42–49}
\begin{leancode}
\leanline{0}{theorem\ IsSubgradient.monotone\ \{f\ :\ E\ →\ ℝ\}\ \{p\ q\ g\ h\ :\ E\}}
\leanline{4}{(hp\ :\ IsSubgradient\ f\ p\ g)\ (hq\ :\ IsSubgradient\ f\ q\ h)\ :}
\leanline{4}{0\ ≤\ {\leanfallback ⟪}g\ -\ h,\ p\ -\ q{\leanfallback ⟫}}
\end{leancode}
\leanlabel{How the Lean proof follows the argument}
\texttt{prox\allowbreak{}Objective f lam x p} is \(Q_{\lambda,x}(p)\), the parameter \(\lambda\) being written \texttt{lam}. \texttt{Is\allowbreak{}Prox\allowbreak{}Point f lam x p} is the proposition \texttt{\ensuremath{\forall} y, prox\allowbreak{}Objective f lam x p \ensuremath{\le} prox\allowbreak{}Objective f lam x y}, and \texttt{Is\allowbreak{}Subgradient f p g} is \texttt{\ensuremath{\forall} y, f p + {\leanfallback ⟪}g, y -\allowbreak{} p{\leanfallback ⟫} \ensuremath{\le} f y}, which is (1) with \(s=g\). Here \texttt{{\leanfallback ⟪}u, v{\leanfallback ⟫}} denotes the real inner product (made available by \texttt{open scoped Real\allowbreak{}Inner\allowbreak{}Product\allowbreak{}Space}), and the ambient space is a type \texttt{E} with instances \texttt{[Normed\allowbreak{}Add\allowbreak{}Comm\allowbreak{}Group E] [Inner\allowbreak{}Product\allowbreak{}Space \ensuremath{\mathbb R} E]}, i.e. a real inner-product space; completeness is not assumed. The definitions accept every real \texttt{lam}: Lean's convention \(a/0=0\) makes \texttt{prox\allowbreak{}Objective} a total function, but every theorem of this section that uses it assumes \texttt{hlam : 0 < lam}, so only \(\lambda>0\) matters.




In \texttt{Is\allowbreak{}Subgradient.\allowbreak{}monotone} the hypotheses \texttt{hp : Is\allowbreak{}Subgradient f p g} and \texttt{hq : Is\allowbreak{}Subgradient f q h} are the two supports; \texttt{hpq := hp q} and \texttt{hqp := hq p} are the two displayed instances of (1). The rewrite \texttt{show q -\allowbreak{} p = -\allowbreak{}(p -\allowbreak{} q) by abel} followed by \texttt{inner\_\allowbreak{}neg\_\allowbreak{}right} turns \(\langle g,q-p\rangle\) into \(-\langle g,p-q\rangle\); \texttt{simp only [inner\_\allowbreak{}sub\_\allowbreak{}left]} expands the goal \(\langle g-h,p-q\rangle\) into \(\langle g,p-q\rangle-\langle h,p-q\rangle\), and \texttt{linarith} adds the two instances exactly as in the proof.
\end{leanbox}

\section{\texorpdfstring{The residual subgradient condition characterizes proximal points}{The residual subgradient condition characterizes proximal points}}


Recall that \(f:E\to\mathbb R\) is \emph{convex} if \(f((1-t)p+ty)\le(1-t)f(p)+tf(y)\) for all \(p,y\in E\) and all \(t\in[0,1]\).




\textbf{Claim.} Let \(\lambda>0\), \(x,p\in E\) and \(f:E\to\mathbb R\).



\begin{enumerate}


\item (Sufficiency.) If \(\lambda^{-1}(x-p)\in\partial f(p)\), then \(p\) is a proximal point of \(f\) at \(x\). No convexity of \(f\) is needed.

\item (Necessity.) If \(f\) is convex and \(p\) is a proximal point of \(f\) at \(x\), then \(\lambda^{-1}(x-p)\in\partial f(p)\).

\end{enumerate}



Consequently, for convex \(f\) and \(\lambda>0\),


\[p\text{ is a proximal point of }f\text{ at }x\quad\Longleftrightarrow\quad\frac{x-p}{\lambda}\in\partial f(p).\]


\textbf{Proof of sufficiency.} Fix \(y\in E\). Writing \(y-x=(p-x)+(y-p)\) and expanding the squared norm,


\[\|y-x\|^2=\|p-x\|^2+2\langle p-x,y-p\rangle+\|y-p\|^2.\qquad(1)\]


The hypothesis, tested at \(y\), reads \(f(p)+\lambda^{-1}\langle x-p,y-p\rangle\le f(y)\). Since \(\langle x-p,y-p\rangle=-\langle p-x,y-p\rangle\), this is


\[f(y)-f(p)\ge-\frac1\lambda\langle p-x,y-p\rangle.\qquad(2)\]


Using (1) and then (2),


\[\begin{aligned}
Q_{\lambda,x}(y)-Q_{\lambda,x}(p)&=f(y)-f(p)+\frac{\|y-x\|^2-\|p-x\|^2}{2\lambda}\\
&=f(y)-f(p)+\frac1\lambda\langle p-x,y-p\rangle+\frac{\|y-p\|^2}{2\lambda}\\
&\ge\frac{\|y-p\|^2}{2\lambda}\ge0 .
\end{aligned}\]


Hence \(Q_{\lambda,x}(p)\le Q_{\lambda,x}(y)\); as \(y\) was arbitrary, \(p\) is a proximal point. \(\square\) The computation also gives the quadratic growth \(Q_{\lambda,x}(y)\ge Q_{\lambda,x}(p)+\|y-p\|^2/(2\lambda)\).




\textbf{Proof of necessity.} Fix \(y\in E\). It suffices to show that


\[f(p)+\lambda^{-1}\langle x-p,y-p\rangle\le f(y)+\varepsilon\qquad\text{for every }\varepsilon>0,\]


because a real number that is at most \(f(y)+\varepsilon\) for every \(\varepsilon>0\) is at most \(f(y)\). Fix \(\varepsilon>0\) and put


\[t=\min\Bigl(1,\ \frac{2\lambda\varepsilon}{\|y-p\|^2+1}\Bigr).\]


Then \(0<t\le1\) and \(t(\|y-p\|^2+1)\le2\lambda\varepsilon\); since \(t>0\), in particular


\[t\|y-p\|^2\le2\lambda\varepsilon.\qquad(3)\]


Let \(y_t=(1-t)p+ty\). Convexity of \(f\) with weights \(1-t\ge0\) and \(t\ge0\) gives


\[f(y_t)\le(1-t)f(p)+tf(y).\qquad(4)\]


Since \(y_t-x=(p-x)+t(y-p)\), expanding the squared norm gives \(\|y_t-x\|^2=\|p-x\|^2+2t\langle p-x,y-p\rangle+t^2\|y-p\|^2\). Minimality of \(p\) gives \(Q_{\lambda,x}(p)\le Q_{\lambda,x}(y_t)\); cancelling \(\|p-x\|^2/(2\lambda)\) and applying (4),


\[f(p)\le f(y_t)+\frac{2t\langle p-x,y-p\rangle+t^2\|y-p\|^2}{2\lambda}\le(1-t)f(p)+tf(y)+\frac{t}{\lambda}\langle p-x,y-p\rangle+\frac{t^2}{2\lambda}\|y-p\|^2.\]


Subtracting \((1-t)f(p)\) from both sides and multiplying by \(2\lambda>0\) gives


\[t\Bigl(2\lambda\bigl(f(p)-f(y)\bigr)-2\langle p-x,y-p\rangle-t\|y-p\|^2\Bigr)\le0.\]


Because \(t>0\), the bracket is nonpositive. Dividing the bracket by \(2\lambda>0\) and using (3) divided by \(2\lambda\),


\[f(p)-\frac1\lambda\langle p-x,y-p\rangle\le f(y)+\frac{t\|y-p\|^2}{2\lambda}\le f(y)+\varepsilon.\]


The left-hand side equals \(f(p)+\lambda^{-1}\langle x-p,y-p\rangle\), which proves the reduced claim. \(\square\)




The equivalence is the conjunction of the two implications. The choice of \(t\) replaces the limit \(t\downarrow0\) of the directional-derivative argument by a single admissible value of \(t\) for each \(\varepsilon\).

\begin{leanbox}{Lean: \texttt{isProxPoint\_\allowbreak{}of\_\allowbreak{}subgradient}, \texttt{subgradient\_\allowbreak{}of\_\allowbreak{}isProxPoint}, \texttt{isProxPoint\_\allowbreak{}iff\_\allowbreak{}subgradient}}
\leansource{\texttt{isProxPoint\_\allowbreak{}of\_\allowbreak{}subgradient} \quad ProximalLangevin/\allowbreak{}Proximal.lean, lines 52–71}
\begin{leancode}
\leanline{0}{theorem\ isProxPoint\_of\_subgradient\ \{f\ :\ E\ →\ ℝ\}\ \{lam\ :\ ℝ\}\ (hlam\ :\ 0\ <\ lam)}
\leanline{4}{\{x\ p\ :\ E\}\ (hp\ :\ IsSubgradient\ f\ p\ (lam⁻¹\ •\ (x\ -\ p)))\ :}
\leanline{4}{IsProxPoint\ f\ lam\ x\ p}
\end{leancode}
\leansource{\texttt{subgradient\_\allowbreak{}of\_\allowbreak{}isProxPoint} \quad ProximalLangevin/\allowbreak{}Proximal.lean, lines 74–113}
\begin{leancode}
\leanline{0}{theorem\ subgradient\_of\_isProxPoint\ \{f\ :\ E\ →\ ℝ\}\ (hf\ :\ ConvexOn\ ℝ\ univ\ f)}
\leanline{4}{\{lam\ :\ ℝ\}\ (hlam\ :\ 0\ <\ lam)\ \{x\ p\ :\ E\}\ (hp\ :\ IsProxPoint\ f\ lam\ x\ p)\ :}
\leanline{4}{IsSubgradient\ f\ p\ (lam⁻¹\ •\ (x\ -\ p))}
\end{leancode}
\leansource{\texttt{isProxPoint\_\allowbreak{}iff\_\allowbreak{}subgradient} \quad ProximalLangevin/\allowbreak{}Proximal.lean, lines 116–119}
\begin{leancode}
\leanline{0}{theorem\ isProxPoint\_iff\_subgradient\ \{f\ :\ E\ →\ ℝ\}\ (hf\ :\ ConvexOn\ ℝ\ univ\ f)}
\leanline{4}{\{lam\ :\ ℝ\}\ (hlam\ :\ 0\ <\ lam)\ \{x\ p\ :\ E\}\ :}
\leanline{4}{IsProxPoint\ f\ lam\ x\ p\ ↔\ IsSubgradient\ f\ p\ (lam⁻¹\ •\ (x\ -\ p))}
\end{leancode}
\leanlabel{How the Lean proof follows the argument}
\texttt{is\allowbreak{}Prox\allowbreak{}Point\_\allowbreak{}of\_\allowbreak{}subgradient} assumes \texttt{hlam : 0 < lam} and \texttt{hp : Is\allowbreak{}Subgradient f p (lam\ensuremath{{}^{-1}} \ensuremath{\bullet} (x -\allowbreak{} p))}, where \texttt{lam\ensuremath{{}^{-1}} \ensuremath{\bullet} (x -\allowbreak{} p)} is the scalar multiple \(\lambda^{-1}(x-p)\); there is no convexity hypothesis. After \texttt{intro y}, \texttt{h := hp y} is the support inequality at \(y\), and \texttt{real\_\allowbreak{}inner\_\allowbreak{}smul\_\allowbreak{}left} moves the scalar \(\lambda^{-1}\) out of the inner product. \texttt{heq} is the expansion (1), obtained from \texttt{norm\_\allowbreak{}add\_\allowbreak{}sq\_\allowbreak{}real} (\(\|u+v\|^2=\|u\|^2+2\langle u,v\rangle+\|v\|^2\)), and \texttt{hi} is the sign change that turns the hypothesis into (2). The goal is multiplied by \(2\lambda>0\) with \texttt{mul\_\allowbreak{}le\_\allowbreak{}mul\_\allowbreak{}iff\_\allowbreak{}of\_\allowbreak{}pos\_\allowbreak{}right}; \texttt{h'} is (2) multiplied by \(\lambda\), and \texttt{field\_\allowbreak{}simp} followed by \texttt{nlinarith [sq\_\allowbreak{}nonneg \ensuremath{\Vert}y -\allowbreak{} p\ensuremath{\Vert}]} performs the displayed computation, the extra term being \(\|y-p\|^2\ge0\).




\texttt{subgradient\_\allowbreak{}of\_\allowbreak{}is\allowbreak{}Prox\allowbreak{}Point} assumes \texttt{hf : Convex\allowbreak{}On \ensuremath{\mathbb R} univ f} (convexity of \(f\) on the whole space \texttt{univ}), \texttt{hlam} and \texttt{hp : Is\allowbreak{}Prox\allowbreak{}Point f lam x p}.



\begin{enumerate}


\item \texttt{apply le\_\allowbreak{}of\_\allowbreak{}forall\_\allowbreak{}pos\_\allowbreak{}le\_\allowbreak{}add} reduces the goal to the inequality with \(+\varepsilon\) for every \texttt{eps} \(>0\).

\item \texttt{t} is defined by the same minimum; \texttt{ht}, \texttt{ht1} and \texttt{hteps} are \(t>0\), \(t\le1\) and (3), the last one from \texttt{le\_\allowbreak{}div\_\allowbreak{}iff\ensuremath{{}_{0}}} applied to \texttt{min\_\allowbreak{}le\_\allowbreak{}right}.

\item \texttt{hconv} is (4): it is the second field \texttt{hf.\allowbreak{}2} of the convexity hypothesis, applied with weights \(1-t\) and \(t\). \texttt{hmin} is minimality of \(p\) tested at \(y_t\), and \texttt{heq} is \(y_t-x=(p-x)+t(y-p)\); the rewrites \texttt{norm\_\allowbreak{}add\_\allowbreak{}sq\_\allowbreak{}real}, \texttt{real\_\allowbreak{}inner\_\allowbreak{}smul\_\allowbreak{}right}, \texttt{norm\_\allowbreak{}smul} and \texttt{abs\_\allowbreak{}of\_\allowbreak{}pos ht} expand \(\|y_t-x\|^2\).

\item \texttt{hstep} is the inequality \(f(p)-\lambda^{-1}\langle p-x,y-p\rangle\le f(y)+t\|y-p\|^2/(2\lambda)\). Its inner \texttt{have} is the product inequality with the factor \(t\); \texttt{nonpos\_\allowbreak{}of\_\allowbreak{}mul\_\allowbreak{}nonpos\_\allowbreak{}right} (from \(ab\le0\) and \(a>0\) infer \(b\le0\)) removes the factor \(t\), giving \texttt{hcancel}.

\item \texttt{herr} is (3) divided by \(2\lambda\), and \texttt{linarith} combines it with \texttt{hstep}.

\end{enumerate}



\texttt{is\allowbreak{}Prox\allowbreak{}Point\_\allowbreak{}iff\_\allowbreak{}subgradient} packages the two implications as an equivalence \texttt{↔}, using convexity only for the forward direction.
\end{leanbox}

\section{\texorpdfstring{Firm nonexpansiveness and uniqueness of proximal points}{Firm nonexpansiveness and uniqueness of proximal points}}


\textbf{Claim.} Let \(f\) be convex, \(\lambda>0\), let \(p\) be a proximal point of \(f\) at \(x\) and \(q\) a proximal point of \(f\) at \(y\), both with parameter \(\lambda\). Then


\[\|p-q\|^2\le\langle x-y,p-q\rangle,\qquad(1)\]
\[\|p-q\|\le\|x-y\|,\qquad\|(x-p)-(y-q)\|\le\|x-y\|.\qquad(2)\]


If \(x=y\), then \(p=q\); that is, the proximal point at a given input is unique.




\textbf{Proof.} By the necessity part of the previous step, \(\lambda^{-1}(x-p)\in\partial f(p)\) and \(\lambda^{-1}(y-q)\in\partial f(q)\). Monotonicity of subgradients gives


\[0\le\bigl\langle\lambda^{-1}(x-p)-\lambda^{-1}(y-q),\,p-q\bigr\rangle=\lambda^{-1}\bigl\langle(x-p)-(y-q),\,p-q\bigr\rangle.\]


Since \(\lambda^{-1}>0\), \(\langle(x-p)-(y-q),p-q\rangle\ge0\). Writing \((x-p)-(y-q)=(x-y)-(p-q)\) and using \(\langle p-q,p-q\rangle=\|p-q\|^2\), this is \(\langle x-y,p-q\rangle-\|p-q\|^2\ge0\), which is (1).




\emph{Uniqueness.} If \(x=y\), the right-hand side of (1) is \(\langle0,p-q\rangle=0\). Hence \(\|p-q\|^2\le0\), so \(\|p-q\|=0\) and \(p=q\).




\emph{Nonexpansiveness.} By (1) and the Cauchy–Schwarz inequality \(\langle x-y,p-q\rangle\le\|x-y\|\,\|p-q\|\),


\[\|p-q\|^2\le\|x-y\|\,\|p-q\|.\]


If \(\|p-q\|>\|x-y\|\), then \(\|p-q\|>0\) and multiplying the strict inequality by \(\|p-q\|\) would give \(\|p-q\|^2>\|x-y\|\,\|p-q\|\), a contradiction. Hence \(\|p-q\|\le\|x-y\|\).




\emph{Residual nonexpansiveness.} Put \(u=x-y\) and \(v=p-q\), so that \((x-p)-(y-q)=u-v\) and (1) reads \(\|v\|^2\le\langle u,v\rangle\). Expanding the square,


\[\|u-v\|^2=\|u\|^2-2\langle u,v\rangle+\|v\|^2\le\|u\|^2-2\|v\|^2+\|v\|^2=\|u\|^2-\|v\|^2\le\|u\|^2.\]


Both norms are nonnegative, so \(\|u-v\|\le\|u\|\), which is the second inequality in (2). \(\square\)




Inequality (1) is the firm nonexpansiveness of the proximal map. By (2), the proximal map \(x\mapsto p\) and the residual map \(x\mapsto x-p\) are both \(1\)-Lipschitz.

\begin{leanbox}{Lean: \texttt{IsProxPoint.\allowbreak{}firm}, \texttt{IsProxPoint.\allowbreak{}unique}, \texttt{IsProxPoint.\allowbreak{}nonexpansive} and 1 more}
\leansource{\texttt{IsProxPoint.\allowbreak{}firm} \quad ProximalLangevin/\allowbreak{}Proximal.lean, lines 122–133}
\begin{leancode}
\leanline{0}{theorem\ IsProxPoint.firm\ \{f\ :\ E\ →\ ℝ\}\ (hf\ :\ ConvexOn\ ℝ\ univ\ f)}
\leanline{4}{\{lam\ :\ ℝ\}\ (hlam\ :\ 0\ <\ lam)\ \{x\ y\ p\ q\ :\ E\}}
\leanline{4}{(hp\ :\ IsProxPoint\ f\ lam\ x\ p)\ (hq\ :\ IsProxPoint\ f\ lam\ y\ q)\ :}
\leanline{4}{‖p\ -\ q‖\ \textasciicircum{}\ 2\ ≤\ {\leanfallback ⟪}x\ -\ y,\ p\ -\ q{\leanfallback ⟫}}
\end{leancode}
\leansource{\texttt{IsProxPoint.\allowbreak{}unique} \quad ProximalLangevin/\allowbreak{}Proximal.lean, lines 136–142}
\begin{leancode}
\leanline{0}{theorem\ IsProxPoint.unique\ \{f\ :\ E\ →\ ℝ\}\ (hf\ :\ ConvexOn\ ℝ\ univ\ f)}
\leanline{4}{\{lam\ :\ ℝ\}\ (hlam\ :\ 0\ <\ lam)\ \{x\ p\ q\ :\ E\}}
\leanline{4}{(hp\ :\ IsProxPoint\ f\ lam\ x\ p)\ (hq\ :\ IsProxPoint\ f\ lam\ x\ q)\ :\ p\ =\ q}
\end{leancode}
\leansource{\texttt{IsProxPoint.\allowbreak{}nonexpansive} \quad ProximalLangevin/\allowbreak{}Proximal.lean, lines 145–151}
\begin{leancode}
\leanline{0}{theorem\ IsProxPoint.nonexpansive\ \{f\ :\ E\ →\ ℝ\}\ (hf\ :\ ConvexOn\ ℝ\ univ\ f)}
\leanline{4}{\{lam\ :\ ℝ\}\ (hlam\ :\ 0\ <\ lam)\ \{x\ y\ p\ q\ :\ E\}}
\leanline{4}{(hp\ :\ IsProxPoint\ f\ lam\ x\ p)\ (hq\ :\ IsProxPoint\ f\ lam\ y\ q)\ :}
\leanline{4}{‖p\ -\ q‖\ ≤\ ‖x\ -\ y‖}
\end{leancode}
\leansource{\texttt{IsProxPoint.\allowbreak{}residual\_\allowbreak{}nonexpansive} \quad ProximalLangevin/\allowbreak{}Proximal.lean, lines 154–162}
\begin{leancode}
\leanline{0}{theorem\ IsProxPoint.residual\_nonexpansive\ \{f\ :\ E\ →\ ℝ\}\ (hf\ :\ ConvexOn\ ℝ\ univ\ f)}
\leanline{4}{\{lam\ :\ ℝ\}\ (hlam\ :\ 0\ <\ lam)\ \{x\ y\ p\ q\ :\ E\}}
\leanline{4}{(hp\ :\ IsProxPoint\ f\ lam\ x\ p)\ (hq\ :\ IsProxPoint\ f\ lam\ y\ q)\ :}
\leanline{4}{‖(x\ -\ p)\ -\ (y\ -\ q)‖\ ≤\ ‖x\ -\ y‖}
\end{leancode}
\leanlabel{How the Lean proof follows the argument}
The four theorems assume \texttt{hf : Convex\allowbreak{}On \ensuremath{\mathbb R} univ f} and \texttt{hlam : 0 < lam} and take proximal points \texttt{hp : Is\allowbreak{}Prox\allowbreak{}Point f lam x p} and \texttt{hq : Is\allowbreak{}Prox\allowbreak{}Point f lam y q}; in \texttt{Is\allowbreak{}Prox\allowbreak{}Point.\allowbreak{}unique} both are at the same input \texttt{x}. The dot notation \texttt{hp.\allowbreak{}firm hf hlam hq} applies \texttt{Is\allowbreak{}Prox\allowbreak{}Point.\allowbreak{}firm} to \texttt{hp}.



\begin{enumerate}


\item \texttt{Is\allowbreak{}Prox\allowbreak{}Point.\allowbreak{}firm} is (1). Its \texttt{h} applies \texttt{Is\allowbreak{}Subgradient.\allowbreak{}monotone} to the two residual subgradients produced by \texttt{subgradient\_\allowbreak{}of\_\allowbreak{}is\allowbreak{}Prox\allowbreak{}Point}. \texttt{rw [← smul\_\allowbreak{}sub, real\_\allowbreak{}inner\_\allowbreak{}smul\_\allowbreak{}left] at h} factors \(\lambda^{-1}\) out of the inner product, and \texttt{nonneg\_\allowbreak{}of\_\allowbreak{}mul\_\allowbreak{}nonneg\_\allowbreak{}right h (inv\_\allowbreak{}pos.\allowbreak{}mpr hlam)} removes the positive factor (the lemma states: \(0\le ab\) and \(0<a\) imply \(0\le b\)). \texttt{heq} is the regrouping \((x-p)-(y-q)=(x-y)-(p-q)\); \texttt{inner\_\allowbreak{}sub\_\allowbreak{}left} and \texttt{real\_\allowbreak{}inner\_\allowbreak{}self\_\allowbreak{}eq\_\allowbreak{}norm\_\allowbreak{}sq} rewrite \(\langle (x-y)-(p-q),p-q\rangle\) as \(\langle x-y,p-q\rangle-\|p-q\|^2\), and \texttt{linarith} concludes.

\item \texttt{Is\allowbreak{}Prox\allowbreak{}Point.\allowbreak{}unique} applies (1) with equal inputs; \texttt{simp only [sub\_\allowbreak{}self, inner\_\allowbreak{}zero\_\allowbreak{}left]} turns the right-hand side into \(0\), \texttt{nlinarith} gives \(\|p-q\|=0\), and \texttt{norm\_\allowbreak{}eq\_\allowbreak{}zero} and \texttt{sub\_\allowbreak{}eq\_\allowbreak{}zero} give \(p=q\).

\item \texttt{Is\allowbreak{}Prox\allowbreak{}Point.\allowbreak{}nonexpansive} combines (1) with the Cauchy–Schwarz inequality \texttt{real\_\allowbreak{}inner\_\allowbreak{}le\_\allowbreak{}norm}; \texttt{nlinarith}, supplied with \(\|p-q\|\ge0\) and \(\|x-y\|\ge0\), performs the contradiction argument of the proof.

\item \texttt{Is\allowbreak{}Prox\allowbreak{}Point.\allowbreak{}residual\_\allowbreak{}nonexpansive} rewrites with the same regrouping, expands \(\|u-v\|^2\) with \texttt{norm\_\allowbreak{}sub\_\allowbreak{}sq\_\allowbreak{}real} (\(\|u-v\|^2=\|u\|^2-2\langle u,v\rangle+\|v\|^2\)), and concludes with \texttt{nlinarith} from (1) and nonnegativity of the two norms.

\end{enumerate}
\end{leanbox}

\section{\texorpdfstring{Existence of proximal points and the proximal map}{Existence of proximal points and the proximal map}}


\textbf{Claim.} Let \(E\) be finite-dimensional, let \(f:E\to\mathbb R\) be convex and bounded below, say \(f(z)\ge b\) for all \(z\in E\) and some \(b\in\mathbb R\), and let \(\lambda>0\). Then for every \(x\in E\) there is a proximal point of \(f\) at \(x\). Consequently the \emph{proximal map}


\[P_\lambda(x)=\text{the unique proximal point of }f\text{ at }x\qquad(x\in E)\]


is well defined, and \(Q_{\lambda,x}(P_\lambda(x))\le Q_{\lambda,x}(y)\) for all \(x,y\in E\).




\textbf{Proof.} Fix \(x\in E\). Since \(f(x)\ge b\), the number \(s=\sqrt{2\lambda(f(x)-b)}\) is a well-defined nonnegative real, and \(R=s+1>0\).




\emph{Continuity.} A convex real-valued function on a finite-dimensional real normed space is locally Lipschitz, in particular continuous. The map \(p\mapsto\|p-x\|^2/(2\lambda)\) is continuous as a composition of continuous maps. Hence \(Q_{\lambda,x}\) is continuous.




\emph{A minimizer on a ball.} In a finite-dimensional normed space closed balls are compact (Heine–Borel). The closed ball \(\bar B=\{p\in E:\|p-x\|\le R\}\) is therefore compact, and it is nonempty because it contains \(x\). A continuous real function on a nonempty compact set attains its minimum, so there is \(p\in\bar B\) with \(Q_{\lambda,x}(p)\le Q_{\lambda,x}(y)\) for every \(y\in\bar B\). Taking \(y=x\),


\[Q_{\lambda,x}(p)\le Q_{\lambda,x}(x)=f(x).\qquad(1)\]


\emph{Competitors outside the ball.} Let \(\|y-x\|>R\). Then \(\|y-x\|>s\ge0\), so \(\|y-x\|^2\ge s^2=2\lambda(f(x)-b)\); dividing by \(2\lambda>0\) gives \(f(x)-b\le\|y-x\|^2/(2\lambda)\). Using \(f(y)\ge b\) and then (1),


\[Q_{\lambda,x}(y)=f(y)+\frac{\|y-x\|^2}{2\lambda}\ge b+\bigl(f(x)-b\bigr)=f(x)\ge Q_{\lambda,x}(p).\]


Together with the minimality on \(\bar B\), this gives \(Q_{\lambda,x}(p)\le Q_{\lambda,x}(y)\) for every \(y\in E\), so \(p\) is a proximal point of \(f\) at \(x\).




\emph{The proximal map.} By the uniqueness proved in the previous step, the proximal point at each input is a single element of \(E\), denoted \(P_\lambda(x)\); by definition it minimizes \(Q_{\lambda,x}\). \(\square\)




The lower bound is used only to make the ball radius finite; convexity is used for continuity, for uniqueness and, through the previous step, for all later properties of \(P_\lambda\).

\begin{leanbox}{Lean: \texttt{exists\_\allowbreak{}proxPoint}, \texttt{prox}, \texttt{prox\_\allowbreak{}isProxPoint}}
\leansource{\texttt{exists\_\allowbreak{}proxPoint} \quad ProximalLangevin/\allowbreak{}Proximal.lean, lines 166–195}
\begin{leancode}
\leanline{0}{theorem\ exists\_proxPoint\ [FiniteDimensional\ ℝ\ E]\ \{f\ :\ E\ →\ ℝ\}}
\leanline{4}{(hf\ :\ ConvexOn\ ℝ\ univ\ f)\ (hb\ :\ ∃\ b,\ ∀\ z,\ b\ ≤\ f\ z)}
\leanline{4}{\{lam\ :\ ℝ\}\ (hlam\ :\ 0\ <\ lam)\ (x\ :\ E)\ :\ ∃\ p,\ IsProxPoint\ f\ lam\ x\ p}
\end{leancode}
\leansource{\texttt{prox} \quad ProximalLangevin/\allowbreak{}Proximal.lean, lines 198–200}
\begin{leancode}
\leanline{0}{def\ prox\ (f\ :\ E\ →\ ℝ)\ (lam\ :\ ℝ)\ (x\ :\ E)\ :\ E\ :=\ by}
\leanline{2}{classical}
\leanline{2}{exact\ if\ h\ :\ ∃\ p,\ IsProxPoint\ f\ lam\ x\ p\ then\ Classical.choose\ h\ else\ 0}
\end{leancode}
\leansource{\texttt{prox\_\allowbreak{}isProxPoint} \quad ProximalLangevin/\allowbreak{}Proximal.lean, lines 203–211}
\begin{leancode}
\leanline{0}{theorem\ prox\_isProxPoint\ [FiniteDimensional\ ℝ\ E]\ \{f\ :\ E\ →\ ℝ\}}
\leanline{4}{(hf\ :\ ConvexOn\ ℝ\ univ\ f)\ (hb\ :\ ∃\ b,\ ∀\ z,\ b\ ≤\ f\ z)}
\leanline{4}{\{lam\ :\ ℝ\}\ (hlam\ :\ 0\ <\ lam)\ (x\ :\ E)\ :\ IsProxPoint\ f\ lam\ x\ (prox\ f\ lam\ x)}
\end{leancode}
\leanlabel{How the Lean proof follows the argument}
\texttt{exists\_\allowbreak{}prox\allowbreak{}Point} assumes the instance \texttt{[Finite\allowbreak{}Dimensional \ensuremath{\mathbb R} E]}, \texttt{hf : Convex\allowbreak{}On \ensuremath{\mathbb R} univ f}, the lower bound \texttt{hb : \ensuremath{\exists} b, \ensuremath{\forall} z, b \ensuremath{\le} f z} and \texttt{hlam : 0 < lam}, and concludes \texttt{\ensuremath{\exists} p, Is\allowbreak{}Prox\allowbreak{}Point f lam x p}. The proof follows the paragraphs above.



\begin{enumerate}


\item \texttt{R} is \(\sqrt{2\lambda(f(x)-b)}+1\), written with \texttt{Real.\allowbreak{}sqrt}; \texttt{hxb} is \(f(x)-b\ge0\) and \texttt{h\allowbreak{}R} is \(R>0\).

\item \texttt{hcont} is continuity of the objective. The term \texttt{hf.\allowbreak{}locally\allowbreak{}Lipschitz} is the Mathlib theorem \texttt{Convex\allowbreak{}On.\allowbreak{}locally\allowbreak{}Lipschitz} (a convex function on all of a finite-dimensional space is locally Lipschitz), and the penalty is continuous by \texttt{continuous\_\allowbreak{}id.\allowbreak{}sub continuous\_\allowbreak{}const}, the norm, the square and \texttt{.\allowbreak{}div\_\allowbreak{}const}.

\item \texttt{(is\allowbreak{}Compact\_\allowbreak{}closed\allowbreak{}Ball x R).\allowbreak{}exists\_\allowbreak{}is\allowbreak{}Min\allowbreak{}On} provides a minimizer \texttt{p} of the objective on \texttt{closed\allowbreak{}Ball x R}, which is nonempty by \texttt{mem\_\allowbreak{}closed\allowbreak{}Ball\_\allowbreak{}self}. Compactness of closed balls holds because finite-dimensional real normed spaces are proper.

\item \texttt{by\_\allowbreak{}cases hy : y \ensuremath{\in} closed\allowbreak{}Ball x R} splits the competitors. Inside the ball, \texttt{hmin hy} is the minimality. Outside, \texttt{hpx} is (1), \texttt{hy\allowbreak{}R} is \(R<\|y-x\|\), \texttt{hsqrt} is \(s^2=2\lambda(f(x)-b)\) (from \texttt{Real.\allowbreak{}sq\_\allowbreak{}sqrt}), \texttt{hsq} and \texttt{hquot} are the two displayed consequences, and \texttt{linarith [hb y]} adds the lower bound \(f(y)\ge b\).

\end{enumerate}



The definition \texttt{prox f lam x} is \texttt{if h : \ensuremath{\exists} p, Is\allowbreak{}Prox\allowbreak{}Point f lam x p then Classical.\allowbreak{}choose h else 0}. When a proximal point exists it returns one chosen by the axiom of choice (\texttt{Classical.\allowbreak{}choose}); otherwise it returns the value \(0\). Under the hypotheses of \texttt{exists\_\allowbreak{}prox\allowbreak{}Point} the first branch is taken, and by \texttt{Is\allowbreak{}Prox\allowbreak{}Point.\allowbreak{}unique} the chosen element is the unique proximal point, so \texttt{prox f lam x} is \(P_\lambda(x)\). Every later result that uses \texttt{prox} assumes convexity, a lower bound, finite dimension and \(\lambda>0\), so the fallback value is never used. \texttt{prox\_\allowbreak{}is\allowbreak{}Prox\allowbreak{}Point} states that \texttt{prox f lam x} is a proximal point: \texttt{dite\_\allowbreak{}eq\_\allowbreak{}left h} selects the first branch of the definition, and \texttt{Classical.\allowbreak{}choose\_\allowbreak{}spec h} is the defining property of the chosen element.
\end{leanbox}

\end{document}
